\adj{} computes the sensitivity of a performance measure to model inputs by solving an adjoint problem alongside a \mf{} simulation \citep{langevin2017}, an approach described by \cite{hayek2025}. The approach is non-intrusive: the forward model runs unmodified, and the quantities the adjoint needs are recovered from the simulation as it proceeds.

This document records the technical detail behind extensions made after the original publication. It is supplemental: it assumes the formulation of \cite{hayek2025} and describes only what has been added or corrected. Each chapter is self-contained and can be read on its own.

\section*{Scope}

Chapter~\ref{ch:instantaneous} describes the instantaneous performance measure, a third form alongside the direct and residual forms, which averages a quantity over the time steps at which it is observed rather than accumulating it.

Chapter~\ref{ch:storage} describes how the storage terms couple one time step to the next in the adjoint, and how the specific-storage and specific-yield contributions are selected to match what \mf{} does.

Chapters~\ref{ch:lake} and \ref{ch:sfr} describe the treatment of two advanced packages whose dependent variables are solved outside the groundwater flow matrix. Both use the same device, a bordered matrix, and the two chapters share the notation established in Chapter~\ref{ch:lake}.

\section*{Notation}

The symbols below follow \cite{hayek2025} so that this document and the paper can be read together.

The groundwater flow equations for time step $k$ are written

\begin{equation}
	\label{eqn:forward}
	\left[ \mathbf{A} \right]^{k} \mathbf{h}^{k} = \mathbf{b}^{k},
\end{equation}

\noindent where $\left[ \mathbf{A} \right]^{k}$ is the matrix of coefficients, $\mathbf{h}^{k}$ is the vector of heads (L), and $\mathbf{b}^{k}$ is the vector of constant terms (L$^{3}$T$^{-1}$). Under the Newton-Raphson scheme the Jacobian $\left[ \mathbf{J} \right]^{k}$ takes the place of $\left[ \mathbf{A} \right]^{k}$. The residual is

\begin{equation}
	\label{eqn:residual}
	\mathbf{r}^{k} = \left[ \mathbf{A} \right]^{k} \mathbf{h}^{k}
	- \mathbf{b}^{k} = \mathbf{0}.
\end{equation}

A performance measure $J$ is the sum of the contributions $F^{k}$ of the individual time steps,

\begin{equation}
	\label{eqn:pm}
	J = \sum_{k=1}^{N} F^{k},
\end{equation}

\noindent for a simulation of $N$ time steps. The adjoint states $\bm{\lambda}^{k}$ solve the adjoint system, formed from the transpose of equation~\ref{eqn:forward} and solved backward in time, and the sensitivity of $J$ to an input parameter $\alpha_{j}$ follows from them.

Cells are indexed by $n$ and $m$, so $C_{n,m}$ is the conductance between two cells (L$^{2}$T$^{-1}$) and $\eta_{n}$ is the list of cells connected to cell $n$. A superscript $k$ denotes the current time step and $k-1$ the previous one.

A well withdraws at a specified volumetric rate, so the derivative of equation~\ref{eqn:residual} with respect to that rate is a unit vector at the pumped cell and the sensitivity of $J$ to it is the adjoint state itself. That identity is used throughout as a check, because it means the reported well sensitivity is $\bm{\lambda}$ without post-processing.

The chapters that follow add symbols where they describe a \mf{} formulation directly, and those keep the names used in the \mf{} documentation \citep{langevin2017}: $SC1$ for the primary storage coefficient (L$^{2}$), $S_s$ and $S_y$ for specific storage (L$^{-1}$) and specific yield (dimensionless), and $S_F$ for the fractional saturation of a cell (dimensionless).
